\documentclass[sigconf]{acmart}

%% Rights management
\setcopyright{none}
\acmConference[AI4HPCC '26]{2nd Workshop on AI-Assisted Compilation for HPC}{July 6--9, 2026}{Kyoto, Japan}

\usepackage{booktabs}
\usepackage{subcaption}
\usepackage{xspace}
\usepackage{listings}
\usepackage{xcolor}

\lstset{
  basicstyle=\ttfamily\scriptsize,
  keywordstyle=\color{blue}\bfseries,
  commentstyle=\color{gray},
  stringstyle=\color{red},
  breaklines=true,
  frame=single,
  columns=fullflexible,
  keepspaces=true,
  numbers=left,
  numberstyle=\tiny\color{gray},
  xleftmargin=1.5em,
}

\newcommand{\sysname}{TritonPilot\xspace}
\newcommand{\pet}{PET\xspace}

\begin{document}

\title{Analysis-Guided LLM Code Generation for GPU Kernels:\\Bridging Polyhedral Analysis and Triton Synthesis}

\author{Qin Xiao}
% \authornote{TODO: Add affiliation}
\affiliation{
  \institution{TODO: Institution}
  \city{TODO}
  \country{TODO}
}
\email{TODO@example.com}

% Add co-authors as needed
% \author{Second Author}
% \affiliation{...}

\begin{abstract}
Large language models (LLMs) can generate GPU kernel code, but without
understanding parallelization constraints, they frequently produce
implementations that are incorrect (data races, stale reads) or
inefficient (serial execution on GPU, excessive kernel launches). We
present \sysname, a system that bridges classical polyhedral analysis
with LLM-based Triton kernel generation.  \sysname extracts
parallelization properties, data dependences, and memory access patterns
from C loop nests using the Polyhedral Extraction Tool (PET) and LLVM's
dependence analysis, then encodes these as structured natural-language
guidance in the LLM prompt.  On the 30-kernel Polybench/C suite,
\sysname achieves a 100\% correctness pass rate with a median speedup of
3.06$\times$ over single-threaded C on an NVIDIA RTX~3090.  An ablation
study shows that analysis-guided generation passes 2 additional kernels,
wins 18 out of 28 head-to-head comparisons, and improves median speedup
from 1.17$\times$ to 3.06$\times$.  At 8$\times$ problem sizes, median
speedup rises to 28.97$\times$ with all 30 kernels passing.
Beyond loop benchmarks, \sysname generalizes to real-world application
kernels from the Rodinia benchmark suite---achieving up to
24.9$\times$ speedup on SRAD (speckle-reducing anisotropic diffusion)
and 13.2$\times$ on LUD (LU decomposition), with all 6 kernels
passing correctness on the first attempt.
\end{abstract}

\begin{CCSXML}
<ccs2012>
  <concept>
    <concept_id>10011007.10011006.10011041</concept_id>
    <concept_desc>Software and its engineering~Compilers</concept_desc>
    <concept_significance>500</concept_significance>
  </concept>
  <concept>
    <concept_id>10010147.10010169.10010170</concept_id>
    <concept_desc>Computing methodologies~Parallel computing methodologies</concept_desc>
    <concept_significance>300</concept_significance>
  </concept>
</ccs2012>
\end{CCSXML}

\ccsdesc[500]{Software and its engineering~Compilers}
\ccsdesc[300]{Computing methodologies~Parallel computing methodologies}

\keywords{LLM code generation, GPU kernels, Triton, polyhedral analysis, compiler-guided synthesis}

\maketitle

%% ===================================================================
\section{Introduction}
\label{sec:intro}

GPU programming remains a significant barrier to adopting hardware
accelerators for scientific computing.  Writing efficient GPU kernels
requires expertise in thread-level parallelism, memory coalescing,
shared memory management, and synchronization---skills that domain
scientists often lack.  OpenAI's Triton~\cite{tillet2019triton}
simplifies GPU programming by providing a block-level abstraction, but
generating correct Triton code still requires understanding which loop
dimensions are safe to parallelize, how data dependences constrain
execution order, and when synchronization barriers are needed.

Recent work has demonstrated that large language models can generate
functional GPU code from natural-language descriptions or sequential
implementations~\cite{chen2024chatgpt,nichols2024modeling,
ouyang2024kernelbench}.  However, LLMs lack the ability to perform
precise dependence analysis: they cannot determine whether parallelizing
a particular loop dimension introduces write-after-read (WAR) hazards,
whether array accesses across loop phases create stale-cache reads, or
whether reduction operations require specific accumulation strategies.
As a result, LLM-generated GPU code frequently suffers from:
\begin{enumerate}
  \item \textbf{Correctness failures}: data races from parallelizing
    dimensions with carried dependences;
  \item \textbf{Performance anti-patterns}: serial grid launches
    (\texttt{grid=(1,)}) that waste GPU parallelism, scalar loops inside
    kernels instead of vectorized operations, or excessive Python-side
    kernel launches;
  \item \textbf{Synchronization errors}: missing memory barriers between
    phases in stencil computations, leading to stale L1 cache reads.
\end{enumerate}

We present \sysname, a system that addresses these challenges by
integrating classical compiler analysis into the LLM prompting
pipeline.  \sysname uses the Polyhedral Extraction Tool
(\pet)~\cite{verdoolaege2012polyhedral} and LLVM's dependence
analysis~\cite{llvm} to extract parallelization properties from C loop
nests, then translates these into structured natural-language guidance
that steers the LLM toward correct and efficient Triton
implementations.  The key insight is that \emph{analysis results can be
communicated to LLMs as domain-expert advice}, enabling the model to
make informed parallelization decisions without performing dependence
analysis itself.

Our contributions are:
\begin{itemize}
  \item A pipeline that automatically extracts parallelization
    constraints from C code via polyhedral and LLVM analysis, and
    encodes them as structured LLM prompt guidance
    (\S\ref{sec:analysis});
  \item An iterative refinement loop with error-specific feedback that
    recovers from correctness and performance failures
    (\S\ref{sec:pipeline});
  \item A comprehensive evaluation on 30 Polybench/C kernels showing
    100\% pass rate, 3.06$\times$ median speedup at standard sizes and
    28.97$\times$ at 8$\times$ sizes, with ablation evidence that
    analysis guidance is the primary driver (\S\ref{sec:eval});
  \item Generalization to 6 real-world application kernels from the
    Rodinia benchmark suite, achieving 100\% first-attempt pass rate
    with up to 24.9$\times$ speedup (\S\ref{sec:realworld}).
\end{itemize}


%% ===================================================================
\section{System Overview}
\label{sec:overview}

\begin{figure}[t]
  \centering
  \includegraphics[width=\columnwidth]{figures/pipeline_overview.pdf}
  \caption{%
    \sysname pipeline.  C loop nests are analyzed by \pet{} and LLVM to
    extract parallelization properties.  These are encoded as structured
    guidance in the LLM prompt alongside the source code.  The LLM
    generates Triton code, which is validated against a C reference
    and benchmarked.  Failures trigger error-specific retry prompts.}
  \label{fig:pipeline}
\end{figure}

Figure~\ref{fig:pipeline} shows the \sysname pipeline.  Given a C loop
nest (e.g., from the Polybench/C benchmark suite~\cite{polybench}), the
system proceeds in four stages:

\paragraph{Stage 1: Static Analysis.}
The C kernel is analyzed by \pet{} to extract a polyhedral
representation, from which we compute: (a)~parallelizable dimensions
with validity constraints, (b)~WAR dependences with loop-level scoping,
(c)~scalar expansion opportunities, and (d)~reduction operation types.
When \pet{} cannot handle a construct (e.g., non-affine subscripts), we
fall back to LLVM's \texttt{DependenceAnalysis} pass, parsing direction
vectors from the IR to determine which loop levels carry dependences.

\paragraph{Stage 2: Prompt Construction.}
Analysis results are translated into structured natural-language sections
appended to the LLM prompt.  Each section follows a consistent format:
the analysis finding, its implication for GPU parallelization, and a
concrete code-level recommendation.  The prompt also includes the
original C source, array read/write modes, the exact function signature
the generated code must match, and Triton-specific compilation rules.

\paragraph{Stage 3: LLM Generation and Validation.}
The prompt is sent to a Claude Sonnet~4 model~\cite{anthropic2025claude},
which generates a Triton kernel implementation.  The generated code is
tested against a C reference compiled as a shared library: both
implementations are executed on the same random inputs, and outputs are
compared element-wise with kernel-specific tolerances.

\paragraph{Stage 4: Iterative Refinement.}
If validation fails, the system classifies the error (numerical
mismatch, compilation error, runtime exception, or timeout) and
constructs a targeted retry prompt.  For numerical errors, the feedback
includes the maximum absolute error and suggestions for index-bounds
checking.  For performance issues (speedup~$<0.1\times$), the feedback
identifies parallelization anti-patterns.  Up to 10 attempts are made,
with a context reset at attempt~6 to escape local minima.


%% ===================================================================
\section{Analysis-Guided Prompt Engineering}
\label{sec:analysis}

The core technical contribution is the translation of compiler analysis
results into LLM-comprehensible guidance.  We describe the five
principal analysis categories and how each maps to prompt guidance.

\subsection{Parallelization Dimension Analysis}

For each loop nest, \pet{} computes which dimensions can be safely
parallelized by checking whether the iteration-to-memory mapping is
injective (no write conflicts) and whether dependences are only carried
by outer sequential loops.  Each dimension is classified as
\emph{valid} (safe to parallelize), \emph{invalid} with a specific
reason (cross-phase write conflict, sequential context, forward
substitution), or \emph{reduction} (parallelizable with atomic or
tree reduction).

The prompt presents valid options as:
\begin{lstlisting}[language={},frame=none,numbers=none,xleftmargin=0em]
Option 1: Parallelize dim `j` (innermost)
  - Grid: triton.cdiv(NY, BLOCK)
  - Each CTA processes BLOCK elements of j
  - Issue: None (fully parallel)
\end{lstlisting}
Invalid options are shown with explanations to prevent the LLM from
attempting unsafe parallelizations.

\subsection{Write-After-Read Dependence Analysis}

WAR dependences determine whether arrays must be cloned before in-place
update.  \pet{} computes exact WAR sets via polyhedral dependence
analysis; when \pet{} fails, we parse LLVM direction vectors and map
them to loop levels.  The guidance distinguishes three cases:

\begin{enumerate}
  \item \textbf{No WAR}: arrays can be updated in-place.
  \item \textbf{WAR with parallel resolution}: clone the array
    (\texttt{B = A.clone()}) and read from the clone while writing to
    the original.
  \item \textbf{WAR requiring sequential execution}: when all
    dimensions carry dependences (e.g., triangular solvers), the
    guidance recommends \texttt{grid=(1,)} with explicit sequential
    loops inside the kernel.
\end{enumerate}

\subsection{Multi-Phase Decomposition}

Many Polybench kernels contain multiple computational phases with
inter-phase dependences (e.g., \texttt{correlation}: mean computation
$\to$ normalization $\to$ correlation matrix).  The analysis detects
phases via cross-phase write conflicts and recommends splitting into
separate kernel launches connected by a Python driver loop.  This is
critical: LLMs frequently attempt to fuse all phases into a single
kernel, causing incorrect results or serial execution.

\subsection{Reduction and Memory Access Patterns}

The analysis identifies reduction operations (sum, dot product, min/max)
and recommends appropriate Triton primitives (\texttt{tl.sum()},
\texttt{tl.dot()}).  For kernels with opposing reductions (e.g.,
\texttt{bicg}: row reduction on $A$ and column reduction on $A^T$), the
guidance recommends fusing both reductions with row-major iteration for
coalesced memory access.

\subsection{Stencil Synchronization}

For stencil computations (e.g., \texttt{jacobi\_1d}), the analysis
detects cross-phase data dependences within a single CTA.  Triton's
vectorized stores go to L1 cache, and subsequent loads from overlapping
addresses may read stale data.  The guidance inserts
\texttt{tl.debug\_barrier()} calls between phases and includes a
pre-validation check that rejects generated code missing these barriers
before executing the expensive GPU test.


%% ===================================================================
\section{Implementation}
\label{sec:pipeline}

\sysname is implemented in approximately 14,000 lines of Python.  The
analysis frontend comprises 18 modules built on \pet{} (which uses
isl~\cite{verdoolaege2010isl} for polyhedral operations) and LLVM~17
for fallback dependence analysis.

\paragraph{C Reference Infrastructure.}
Each Polybench kernel is compiled as a shared library (\texttt{.so})
with \texttt{-O0} to preserve the reference semantics.  Global arrays
are accessed via \texttt{ctypes.in\_dll()}, and the test harness
initializes identical \texttt{torch.Tensor} inputs for both the C
reference and the Triton implementation.  Correctness is verified with
element-wise comparison using per-kernel tolerances that account for
FP32 accumulation order differences (e.g., \texttt{atol=6.0} for triple
matrix multiply due to \texttt{tl.dot()} vs.\ sequential accumulation).

\paragraph{Retry Strategy.}
The retry loop makes up to 10 attempts per kernel.  Attempts 1--5 build
on prior errors; attempt~6 resets context (discards failed code to
escape fixation on a broken approach); attempts 7--10 continue from the
reset.  Error feedback is category-specific:
\begin{itemize}
  \item \emph{Numerical}: reports max absolute/relative error, suggests
    loop-bound and indexing checks.
  \item \emph{Missing barrier}: provides exact placement template for
    \texttt{tl.debug\_barrier()} between stencil phases.
  \item \emph{Low speedup} ($<0.1\times$): identifies anti-patterns
    (serial grid, scalar loops, excessive launches).
  \item \emph{Compilation/runtime}: passes error message verbatim.
\end{itemize}

\paragraph{Benchmark Methodology.}
Performance is measured as the ratio of C reference wall-clock time to
Triton kernel time, each averaged over 50 iterations after 5 warmup
runs.  The C reference is compiled with \texttt{-O2} and executed
single-threaded to represent the baseline a domain scientist would
write without GPU expertise.


%% ===================================================================
\section{Evaluation}
\label{sec:eval}

We evaluate \sysname on the 30-kernel Polybench/C~4.2.1 suite on a
system with two NVIDIA RTX~3090 GPUs (24\,GB VRAM each), PyTorch~2.5.1,
Triton~3.1.0, and CUDA~12.4.  The LLM is Claude Sonnet~4
(\texttt{claude-sonnet-4-20250514}).

\subsection{Correctness and Performance (Standard Sizes)}

\begin{table}[t]
  \caption{Polybench/C results at standard (1$\times$) problem sizes.
    WA = with analysis, NA = no analysis.}
  \label{tab:results_1x}
  \centering\small
  \begin{tabular}{lcc}
    \toprule
    \textbf{Metric} & \textbf{WA} & \textbf{NA} \\
    \midrule
    Pass rate              & 30/30 (100\%) & 28/30 (93\%) \\
    First-try pass         & 24/30         & ---           \\
    Median speedup         & 3.06$\times$  & 1.17$\times$ \\
    Mean speedup           & 3.83$\times$  & 2.94$\times$ \\
    Kernels $>1\times$     & 21/30         & 14/28         \\
    \midrule
    \multicolumn{3}{l}{\emph{Head-to-head (28 kernels both passed):}} \\
    WA wins ($>$15\% faster) & \multicolumn{2}{c}{18} \\
    NA wins ($>$15\% faster) & \multicolumn{2}{c}{4}  \\
    Ties (within 15\%)       & \multicolumn{2}{c}{6}  \\
    WA-only passes           & \multicolumn{2}{c}{2 (jacobi\_2d, seidel\_2d)}  \\
    \bottomrule
  \end{tabular}
\end{table}

Table~\ref{tab:results_1x} summarizes the results.  With analysis
guidance, all 30 kernels pass correctness, 24 on the first attempt.
The median speedup of 3.06$\times$ over single-threaded C reflects the
mix of compute-bound kernels that benefit from GPU parallelism (e.g.,
\texttt{heat\_3d}: 12.36$\times$, \texttt{3mm}: 13.04$\times$) and
memory-bound or inherently sequential kernels where Triton overhead
dominates (e.g., \texttt{trisolv}: 1.09$\times$, \texttt{durbin}:
0.45$\times$).

Without analysis, two kernels (\texttt{jacobi\_2d},
\texttt{seidel\_2d}) fail entirely---the LLM parallelizes a dimension
with carried dependences, producing incorrect results across all
retry attempts.  In head-to-head comparison, analysis-guided generation
wins 18 of 28 benchmarked kernels by more than 15\%.

\subsection{Scaling to Realistic Problem Sizes}

\begin{table}[t]
  \caption{Results at 8$\times$ problem sizes (e.g., $N$: 60$\to$480,
    matrix dims: 120$\to$960).  Only benchmarked kernels shown.}
  \label{tab:results_8x}
  \centering\small
  \begin{tabular}{lcc}
    \toprule
    \textbf{Metric} & \textbf{1$\times$} & \textbf{8$\times$} \\
    \midrule
    Pass rate          & 30/30     & 30/30     \\
    Benchmarked        & 30        & 26        \\
    Median speedup     & 3.06$\times$ & 28.97$\times$ \\
    Kernels $>1\times$ & 21/30     & 22/26     \\
    Max speedup        & 13.04$\times$ & 1826$\times$ \\
    \bottomrule
  \end{tabular}
\end{table}

At 8$\times$ problem sizes (Table~\ref{tab:results_8x}), median speedup
increases from 3.06$\times$ to 28.97$\times$ as larger matrices
amortize GPU launch overhead and better utilize the GPU's parallelism.
Four kernels (\texttt{adi}, \texttt{cholesky}, \texttt{nussinov},
\texttt{seidel\_2d}) pass correctness but their Triton implementations
are too slow to benchmark within the 60-second timeout---these are
inherently sequential algorithms where \texttt{grid=(1,)} execution
cannot exploit GPU parallelism.

The top performers at 8$\times$ are matrix-multiply kernels that
benefit from \texttt{tl.dot()} tiled accumulation:
\texttt{3mm} (1826$\times$), \texttt{2mm} (1515$\times$),
\texttt{doitgen} (1283$\times$).


\subsection{Iterative Guidance Refinement}

\begin{table}[t]
  \caption{Impact of five guidance categories.  ``Before'' and
    ``After'' are speedups at 1$\times$ problem sizes except where noted.}
  \label{tab:guidance_fixes}
  \centering\small
  \begin{tabular}{llcc}
    \toprule
    \textbf{Guidance Category} & \textbf{Kernel} & \textbf{Before} & \textbf{After} \\
    \midrule
    Multi-phase split      & correlation & 0.19$\times$ & 2.50$\times$ \\
    WAR sequential fallback & nussinov   & FAIL          & PASS         \\
    Triangular grid        & lud$^\dagger$ & 0.14$\times$ & 13.18$\times$ \\
    Fused reductions       & bicg        & 1.04$\times$ & 3.51$\times$ \\
    Stencil barriers       & jacobi\_1d  & 0.13$\times$ & 3.72$\times$ \\
    \midrule
    2D stencil threshold$^\ddagger$ & jacobi\_2d & 1.70$\times$ & 36.82$\times$ \\
    2D stencil threshold$^\ddagger$ & fdtd\_2d   & 1.55$\times$ & 15.43$\times$ \\
    \bottomrule
    \multicolumn{4}{l}{\scriptsize $^\dagger$Rodinia benchmark.
      $^\ddagger$At 8$\times$ sizes.}
  \end{tabular}
\end{table}

Table~\ref{tab:guidance_fixes} shows the impact of individual guidance
categories developed through iterative analysis of failure modes.
Each category is triggered by analysis properties (e.g., cross-phase
dependences, triangular access patterns), not per-kernel heuristics.
The largest improvements come from multi-phase decomposition
(correlation: 13$\times$ improvement) and the triangular grid
recommendation (lud: 94$\times$ improvement).

At 8$\times$ sizes, a threshold adjustment for 2D stencils
(recommending multi-CTA grids instead of \texttt{grid=(1,)} when
$\geq2$ valid parallelization dimensions exist) improved
\texttt{jacobi\_2d} from 1.70$\times$ to 36.82$\times$ and
\texttt{fdtd\_2d} from 1.55$\times$ to 15.43$\times$.


\subsection{Analysis of Failures and Limitations}

Four kernels consistently achieve $<1\times$ speedup at both scales:

\begin{itemize}
  \item \textbf{durbin} (0.37$\times$): inherently sequential
    120-step recurrence with loop-carried dependences on every
    dimension.
  \item \textbf{gramschmidt} (0.41$\times$): Modified Gram-Schmidt
    requires sequential column processing; only the inner dot product
    is parallelizable.
  \item \textbf{trisolv} (0.40$\times$): triangular solve with
    forward substitution---O($N^2$) sequential work.
  \item \textbf{ludcmp} (0.91$\times$): LU decomposition with
    forward/backward substitution; partially parallelizable but
    kernel launch overhead dominates.
\end{itemize}

These are \emph{correct} failures---the analysis correctly identifies
these kernels as having limited parallelism, and the LLM generates
working (but slow) GPU implementations.  This is preferable to
incorrect parallelization, which the no-analysis baseline produces
for \texttt{jacobi\_2d} and \texttt{seidel\_2d}.


\subsection{Real-World Application Kernels}
\label{sec:realworld}

To evaluate whether \sysname generalizes beyond loop benchmarks, we
apply it to 6 application kernels from the Rodinia benchmark
suite~\cite{che2009rodinia}: three extracted as self-contained loop
nests (hotspot, lud, pathfinder) and three larger multi-phase kernels
(SRAD, lavaMD, Gaussian elimination).  These kernels exhibit
challenges absent from Polybench: multi-phase computations with
cross-phase dependences, triangular iteration spaces, and irregular
neighbor traversals.

\begin{table}[t]
  \caption{Real-world application kernel results.  All 6 kernels pass
    correctness on the first attempt with analysis guidance.}
  \label{tab:realworld}
  \centering\small
  \begin{tabular}{lcccc}
    \toprule
    \textbf{Kernel} & \textbf{WA} & \textbf{NA} & \textbf{NA} & \textbf{Key} \\
    & \textbf{Speedup} & \textbf{Speedup} & \textbf{Att.} & \textbf{Challenge} \\
    \midrule
    SRAD        & 24.86$\times$ & 3.50$\times$ & 1 & Multi-phase (3 phases) \\
    lud         & 13.18$\times$ & --- & --- & Triangular decomp. \\
    gaussian    & 3.35$\times$  & 3.61$\times$ & 3 & Triangular + back-sub \\
    lavaMD      & 1.47$\times$  & 0.56$\times$ & 5 & Irregular neighbor list \\
    hotspot     & 1.05$\times$  & --- & --- & 2D stencil \\
    pathfinder  & 0.11$\times$  & --- & --- & Sequential DP \\
    \bottomrule
  \end{tabular}
\end{table}

Table~\ref{tab:realworld} shows the results.  All 6 kernels pass
correctness on the first attempt with analysis guidance.  The three
multi-phase kernels (SRAD, lavaMD, Gaussian) also have no-analysis
baselines for comparison:

\paragraph{SRAD (24.9$\times$).}
This speckle-reducing anisotropic diffusion kernel has three
computational phases per iteration (ROI statistics, diffusion
coefficients, divergence update).  The analysis detects cross-phase
dependences and recommends splitting into separate kernels with
parallel grids, yielding a 7.1$\times$ improvement over the
no-analysis baseline (3.5$\times$), which serializes all phases.

\paragraph{Gaussian elimination (3.4$\times$).}
This kernel has a triangular iteration space with cross-phase
dependences between the elimination and back-substitution phases.
The analysis recommends a single kernel with the sequential dimension
as an internal loop, \texttt{tl.sum(mask) > 0} guards to skip
empty-mask CTAs in the triangular region, and back-substitution in
CTA~0.  The no-analysis baseline achieves similar performance
(3.6$\times$) but requires 3 attempts due to initial timeouts from
\texttt{tl.constexpr} misuse.

\paragraph{lavaMD (1.5$\times$).}
This molecular dynamics kernel traverses irregular neighbor lists.
The analysis identifies the particle dimension as parallelizable,
enabling a 2.6$\times$ improvement over the no-analysis baseline
(0.56$\times$), which fails with type-mismatch errors on 4 of 5
attempts.

The real-world results demonstrate two benefits of analysis guidance
beyond loop benchmarks: (1)~correct multi-phase decomposition for
kernels with complex inter-phase dependences (SRAD), and
(2)~first-attempt success vs.\ multiple retries without guidance
(lavaMD, Gaussian).


%% ===================================================================
\section{Related Work}
\label{sec:related}

\paragraph{LLM Code Generation for GPUs.}
Chen et al.~\cite{chen2024chatgpt} evaluate ChatGPT on CUDA kernel
generation from natural-language descriptions, finding frequent
correctness issues.  Nichols et al.~\cite{nichols2024modeling} use LLMs
for performance modeling of GPU kernels.
KernelBench~\cite{ouyang2024kernelbench} benchmarks LLM-generated
Triton/CUDA code on 250 tasks, reporting 20--40\% functional
correctness without guidance.  Our work differs by using compiler
analysis to \emph{guide} the LLM, achieving 100\% correctness on
Polybench/C.

\paragraph{AI-Assisted Compilation.}
CompilerGym~\cite{cummins2022compilergym} provides RL environments for
compiler optimization.  TVM~\cite{chen2018tvm} and
Halide~\cite{ragan2013halide} use auto-scheduling with learned cost
models.  AlphaCode~\cite{li2022alphacode} generates competitive
programming solutions.  \sysname is complementary: rather than
replacing the compiler, it uses compiler analysis to improve LLM-based
code generation.

\paragraph{Polyhedral Compilation.}
Polyhedral frameworks (Pluto~\cite{bondhugula2008pluto},
PPCG~\cite{verdoolaege2013ppcg}, Polly~\cite{grosser2012polly})
automatically generate parallel code from affine loop nests.  These
produce optimal schedules but are limited to affine programs and
produce low-level code.  \sysname uses polyhedral analysis for
\emph{guidance} rather than code generation, enabling the LLM to handle
non-affine constructs, complex data layouts, and high-level Triton
abstractions that polyhedral compilers cannot target directly.


%% ===================================================================
\section{Conclusion}
\label{sec:conclusion}

We presented \sysname, a system that bridges polyhedral compiler
analysis with LLM-based GPU kernel generation.  By translating
parallelization constraints, dependence information, and memory access
patterns into structured prompt guidance, \sysname enables an LLM to
generate correct and efficient Triton kernels for all 30 Polybench/C
benchmarks, with a median speedup of 3.06$\times$ at standard sizes
and 28.97$\times$ at 8$\times$ sizes.  The ablation study confirms that
analysis guidance is essential: without it, 2 kernels fail entirely and
median speedup drops to 1.17$\times$.  Beyond loop benchmarks, the
approach generalizes to 6 real-world Rodinia kernels---all passing on
the first attempt, with SRAD achieving 24.9$\times$ speedup through
analysis-guided multi-phase decomposition.

Future work includes extending the analysis to non-affine programs,
supporting multi-GPU parallelization, and scaling to larger scientific
codebases with more complex data structures.

%% ===================================================================
%% References
\bibliographystyle{ACM-Reference-Format}
\begin{thebibliography}{99}

\bibitem{tillet2019triton}
P.~Tillet, H.~T. Kung, and D.~Cox.
\newblock Triton: An intermediate language and compiler for tiled neural
  network computations.
\newblock In \emph{Proc. MAPL}, 2019.

\bibitem{chen2024chatgpt}
Y.~Chen, S.~Huang, D.~Zheng, et~al.
\newblock {ChatGPT} for {CUDA} programming: Evaluating {LLM}-generated {GPU}
  code.
\newblock \emph{arXiv preprint arXiv:2404.12000}, 2024.

\bibitem{nichols2024modeling}
B.~Nichols, S.~Wang, et~al.
\newblock Modeling {GPU} kernels with {LLMs}.
\newblock \emph{arXiv preprint}, 2024.

\bibitem{ouyang2024kernelbench}
A.~Ouyang, S.~Zheng, et~al.
\newblock {KernelBench}: Can {LLMs} write efficient {GPU} kernels?
\newblock \emph{arXiv preprint arXiv:2404.12489}, 2024.

\bibitem{verdoolaege2012polyhedral}
S.~Verdoolaege and T.~Grosser.
\newblock Polyhedral extraction tool.
\newblock In \emph{Proc. IMPACT}, 2012.

\bibitem{llvm}
C.~Lattner and V.~Adve.
\newblock {LLVM}: A compilation framework for lifelong program analysis \&
  transformation.
\newblock In \emph{Proc. CGO}, 2004.

\bibitem{polybench}
L.-N. Pouchet.
\newblock Polybench/C: The polyhedral benchmark suite, version 4.2.1.
\newblock \url{https://web.cs.ucla.edu/~pouchet/software/polybench/}, 2016.

\bibitem{anthropic2025claude}
Anthropic.
\newblock Claude {Sonnet} 4 model card.
\newblock \url{https://www.anthropic.com/claude}, 2025.

\bibitem{verdoolaege2010isl}
S.~Verdoolaege.
\newblock isl: An integer set library for the polyhedral model.
\newblock In \emph{Proc. ICMS}, 2010.

\bibitem{cummins2022compilergym}
C.~Cummins, B.~Wasti, J.~Guo, et~al.
\newblock {CompilerGym}: Robust, performant compiler optimization
  environments for {AI} research.
\newblock In \emph{Proc. CGO}, 2022.

\bibitem{chen2018tvm}
T.~Chen, T.~Moreau, Z.~Jiang, et~al.
\newblock {TVM}: An automated end-to-end optimizing compiler for deep
  learning.
\newblock In \emph{Proc. OSDI}, 2018.

\bibitem{ragan2013halide}
J.~Ragan-Kelley, C.~Barnes, A.~Adams, et~al.
\newblock Halide: A language and compiler for optimizing parallelism,
  locality, and recomputation in image processing pipelines.
\newblock In \emph{Proc. PLDI}, 2013.

\bibitem{li2022alphacode}
Y.~Li, D.~Choi, J.~Chung, et~al.
\newblock Competition-level code generation with {AlphaCode}.
\newblock \emph{Science}, 378(6624), 2022.

\bibitem{bondhugula2008pluto}
U.~Bondhugula, A.~Hartono, J.~Ramanujam, and P.~Sadayappan.
\newblock A practical automatic polyhedral parallelizer and locality
  optimizer.
\newblock In \emph{Proc. PLDI}, 2008.

\bibitem{verdoolaege2013ppcg}
S.~Verdoolaege, J.~Carlos~Juega, A.~Cohen, et~al.
\newblock Polyhedral parallel code generation for {CUDA}.
\newblock \emph{ACM TACO}, 9(4), 2013.

\bibitem{grosser2012polly}
T.~Grosser, A.~Groesslinger, and C.~Lengauer.
\newblock Polly---performing polyhedral optimizations on a low-level
  intermediate representation.
\newblock \emph{Parallel Processing Letters}, 22(04), 2012.

\bibitem{che2009rodinia}
S.~Che, M.~Boyer, J.~Meng, et~al.
\newblock Rodinia: A benchmark suite for heterogeneous computing.
\newblock In \emph{Proc. IISWC}, 2009.

\end{thebibliography}

\end{document}
